\documentclass[11pt]{article}
\usepackage[margin=2.1cm]{geometry}
\usepackage{booktabs}
\usepackage{amsmath}
\usepackage{enumitem}
\usepackage[T1]{fontenc}
\usepackage{lmodern}

\setlength{\parindent}{0pt}
\setlength{\parskip}{5pt}
\newcommand{\code}[1]{\texttt{#1}}
\newcommand{\answerbox}[1]{\vspace{2pt}\fbox{\parbox[c][#1][c]{\dimexpr\linewidth-2\fboxsep-2\fboxrule}{\ }}\vspace{4pt}}

\begin{document}

\begin{center}
{\Large\bfseries EE\,362 --- Pre-Lab 9}\\[2pt]
{\large Memory: ROM and RAM}\\[6pt]
\begin{tabular}{@{}p{6.5cm}p{4cm}p{4cm}@{}}
Name: \hrulefill & Section: Mon / Tue & Date: \hrulefill
\end{tabular}
\end{center}

\vspace{2pt}\hrule\vspace{4pt}

{\small\textbf{Due at the start of the lab session.} Five points, about twenty
minutes, on paper. Bring the instruction-format sheet.}

\vspace{4pt}\hrule

\section*{Question 1 --- Disassemble the program}

These are the first four words of the program your CPU runs in Lab 12. Assemble
each from its bytes (big-endian), then decode it.

\begin{center}
\begin{tabular}{lll@{\hskip 6pt}||@{\hskip 6pt}l@{\hskip 10pt}l}
\toprule
addr & bytes & & 32-bit word & instruction \\
\midrule
\code{0x00} & \code{3C 01 00 0B} & & \rule{2.8cm}{0.4pt} & \rule{4cm}{0.4pt} \\[5pt]
\code{0x08} & \code{20 21 FF FF} & & \rule{2.8cm}{0.4pt} & \rule{4cm}{0.4pt} \\[5pt]
\code{0x14} & \code{08 00 00 02} & & \rule{2.8cm}{0.4pt} & \rule{4cm}{0.4pt} \\[5pt]
\code{0x18} & \code{AC 02 00 00} & & \rule{2.8cm}{0.4pt} & \rule{4cm}{0.4pt} \\
\bottomrule
\end{tabular}
\end{center}

\textbf{(b)} Take the bytes at \code{0x08} and assemble them the \emph{other}
way round (little-endian). What 32-bit word do you get, and what is its opcode
field? Is that a valid instruction in this ISA?

\answerbox{2.0cm}

\textbf{(c)} In one sentence: why does a program compiled for one byte order
crash on a machine with the other?

\answerbox{1.4cm}

\section*{Question 2 --- The bus}

\textbf{(a)} Complete the RAM control table. \textbf{Careful: \code{rw = '1'} is
READ.}

\begin{center}
\begin{tabular}{cccc@{\hskip 8pt}||@{\hskip 8pt}l@{\hskip 14pt}l}
\toprule
\# & \code{en} & \code{rw} & \code{Strb} & Data bus & memory \\
\midrule
0 & 0 & --- & --- & \rule{2.6cm}{0.4pt} & \rule{3.4cm}{0.4pt} \\[5pt]
1 & 1 & 1 & --- & \rule{2.6cm}{0.4pt} & \rule{3.4cm}{0.4pt} \\[5pt]
2 & 1 & 0 & 0 & \rule{2.6cm}{0.4pt} & \rule{3.4cm}{0.4pt} \\[5pt]
3 & 1 & 0 & 1 & \rule{2.6cm}{0.4pt} & \rule{3.4cm}{0.4pt} \\
\bottomrule
\end{tabular}
\end{center}

\textbf{(b)} Write the two equations:

\vspace{4pt}
DataDrive $=$ \rule{6cm}{0.4pt} \hfill Write $=$ \rule{6cm}{0.4pt}
\vspace{8pt}

\textbf{(c)} During a store, the RAM must \emph{not} drive the data bus even
though it is enabled. Who drives it instead, and what would happen if both did?

\answerbox{2.0cm}

\section*{Question 3 --- Test plan}

\begin{center}
\begin{tabular}{@{}c@{\hskip 8pt}p{5.0cm}@{\hskip 8pt}p{7.6cm}@{}}
\toprule
\# & what you do & what it proves \\
\midrule
1 & \rule{4.8cm}{0.4pt} & \rule{7.4cm}{0.4pt} \\[8pt]
2 & \rule{4.8cm}{0.4pt} & \rule{7.4cm}{0.4pt} \\[8pt]
3 & \rule{4.8cm}{0.4pt} & \rule{7.4cm}{0.4pt} \\[8pt]
4 & \rule{4.8cm}{0.4pt} & \rule{7.4cm}{0.4pt} \\[8pt]
5 & \rule{4.8cm}{0.4pt} & \rule{7.4cm}{0.4pt} \\[8pt]
6 & \rule{4.8cm}{0.4pt} & \rule{7.4cm}{0.4pt} \\
\bottomrule
\end{tabular}
\end{center}

\textbf{Then answer both:}

\textbf{(a)} A classmate tests their memory only by storing \code{0xAABBBBAA}
and reading it back. Why can that test not detect a byte-order error?

\answerbox{1.6cm}

\textbf{(b)} Your ROM has 512 bytes and the address input is 32 bits. What
happens if someone reads address \code{0x000001FF}, and how will you handle it?

\answerbox{2.0cm}

\vfill
{\small\emph{Question 3 is the one that matters.} It is the same question as
Section 5 of your lab report.}

\end{document}
